\documentclass[11pt]{article}
% The preamble every project's paper shares: % The preamble every project's paper shares: % The preamble every project's paper shares: \input{../../preamble} from
% projects/<project>/paper/. Boards come from chessboard and xskak, plots from
% pgfplots; every number a paper prints comes from its generated.tex.
\usepackage[LSBC4,T1]{fontenc}
\usepackage[utf8]{inputenc}
\usepackage{lmodern}
\usepackage[margin=1in]{geometry}
\usepackage{microtype}
\usepackage{amsmath}
\usepackage{booktabs}
\usepackage{tabularx}
\usepackage{array}
\usepackage{xcolor}
\usepackage[obeyspaces,spaces]{url}
\usepackage{xskak}
\usepackage{pgfplots}
\pgfplotsset{compat=1.18}
\usepackage{caption}
\usepackage{subcaption}
\usepackage{enumitem}
\usepackage[section]{placeins}
\usepackage[hidelinks]{hyperref}

\setchessboard{boardfontencoding=LSBC4, boardfontfamily=skaknew, blackfieldcolor=black!22,
  whitefieldcolor=white, setfontcolors, showmover=false, label=true, labelfontsize=6pt,
  boardfontsize=14pt}

% A move or line in UCI, and a position in FEN, both breakable at a slash.
\newcommand{\uci}[1]{\texttt{#1}}
\newcommand{\fen}[1]{{\footnotesize\path{#1}}}
\newcolumntype{L}{>{\raggedright\arraybackslash}X}

\setlength{\parskip}{0.4em}
\setlength{\parindent}{0pt}
 from
% projects/<project>/paper/. Boards come from chessboard and xskak, plots from
% pgfplots; every number a paper prints comes from its generated.tex.
\usepackage[LSBC4,T1]{fontenc}
\usepackage[utf8]{inputenc}
\usepackage{lmodern}
\usepackage[margin=1in]{geometry}
\usepackage{microtype}
\usepackage{amsmath}
\usepackage{booktabs}
\usepackage{tabularx}
\usepackage{array}
\usepackage{xcolor}
\usepackage[obeyspaces,spaces]{url}
\usepackage{xskak}
\usepackage{pgfplots}
\pgfplotsset{compat=1.18}
\usepackage{caption}
\usepackage{subcaption}
\usepackage{enumitem}
\usepackage[section]{placeins}
\usepackage[hidelinks]{hyperref}

\setchessboard{boardfontencoding=LSBC4, boardfontfamily=skaknew, blackfieldcolor=black!22,
  whitefieldcolor=white, setfontcolors, showmover=false, label=true, labelfontsize=6pt,
  boardfontsize=14pt}

% A move or line in UCI, and a position in FEN, both breakable at a slash.
\newcommand{\uci}[1]{\texttt{#1}}
\newcommand{\fen}[1]{{\footnotesize\path{#1}}}
\newcolumntype{L}{>{\raggedright\arraybackslash}X}

\setlength{\parskip}{0.4em}
\setlength{\parindent}{0pt}
 from
% projects/<project>/paper/. Boards come from chessboard and xskak, plots from
% pgfplots; every number a paper prints comes from its generated.tex.
\usepackage[LSBC4,T1]{fontenc}
\usepackage[utf8]{inputenc}
\usepackage{lmodern}
\usepackage[margin=1in]{geometry}
\usepackage{microtype}
\usepackage{amsmath}
\usepackage{booktabs}
\usepackage{tabularx}
\usepackage{array}
\usepackage{xcolor}
\usepackage[obeyspaces,spaces]{url}
\usepackage{xskak}
\usepackage{pgfplots}
\pgfplotsset{compat=1.18}
\usepackage{caption}
\usepackage{subcaption}
\usepackage{enumitem}
\usepackage[section]{placeins}
\usepackage[hidelinks]{hyperref}

\setchessboard{boardfontencoding=LSBC4, boardfontfamily=skaknew, blackfieldcolor=black!22,
  whitefieldcolor=white, setfontcolors, showmover=false, label=true, labelfontsize=6pt,
  boardfontsize=14pt}

% A move or line in UCI, and a position in FEN, both breakable at a slash.
\newcommand{\uci}[1]{\texttt{#1}}
\newcommand{\fen}[1]{{\footnotesize\path{#1}}}
\newcolumntype{L}{>{\raggedright\arraybackslash}X}

\setlength{\parskip}{0.4em}
\setlength{\parindent}{0pt}

% Written by projects/structure-breaks/paper/generate.mjs; do not edit by hand.

\newcommand{\AppCommit}{f4bd63f19b}
\newcommand{\MinMen}{eight}
\newcommand{\ClockFloor}{30}
\newcommand{\WinLow}{10}
\newcommand{\WinHigh}{90}
\newcommand{\MinLegal}{two}
\newcommand{\Offset}{250}
\newcommand{\Above}{200}
\newcommand{\GradeNodes}{250,000}
\newcommand{\Window}{3}
\newcommand{\MinPositions}{150}
\newcommand{\MinCell}{30}
\newcommand{\Draws}{2,000}
\newcommand{\Seed}{20261019}
\newcommand{\StrongTarget}{1,200}
\newcommand{\StrongDrawn}{3,000}
\newcommand{\Tries}{three}
\newcommand{\Games}{647}
\newcommand{\GamesFirst}{324}
\newcommand{\GamesSecond}{323}
\newcommand{\FirstFrom}{July 22}
\newcommand{\FirstTo}{August 23}
\newcommand{\SecondFrom}{August 23}
\newcommand{\SecondTo}{October 15}
\newcommand{\Candidates}{8,896}
\newcommand{\Named}{2,411}
\newcommand{\NotOn}{1,040}
\newcommand{\NoLever}{twelve}
\newcommand{\BreakOn}{1,359}
\newcommand{\BreakOnFirst}{730}
\newcommand{\BreakOnSecond}{629}
\newcommand{\BreakGames}{245}
\newcommand{\BreakGamesFirst}{130}
\newcommand{\BreakGamesSecond}{115}
\newcommand{\NamedShare}{27}
\newcommand{\OnShare}{56}
\newcommand{\LondonOn}{638}
\newcommand{\LondonShare}{47}
\newcommand{\MonthName}{September 2026}
\newcommand{\Floor}{2200}
\newcommand{\StreamCap}{400,000}
\newcommand{\StrongGamesRead}{11,392}
\newcommand{\StrongNoStructure}{5,219}
\newcommand{\StrongNotOn}{3,173}
\newcommand{\Blitz}{2,781}
\newcommand{\Rapid}{209}
\newcommand{\Classical}{10}
\newcommand{\StrongIqp}{1,171}
\newcommand{\StrongIqpShare}{39}
\newcommand{\ScoredFirst}{491}
\newcommand{\ScoredGamesFirst}{114}
\newcommand{\ForcingFirst}{233}
\newcommand{\DecidedFirst}{6}
\newcommand{\NoExactFirst}{0}
\newcommand{\NoOtherFirst}{0}
\newcommand{\InexactFirst}{0}
\newcommand{\ScoredSecond}{418}
\newcommand{\ScoredGamesSecond}{103}
\newcommand{\ForcingSecond}{208}
\newcommand{\DecidedSecond}{3}
\newcommand{\NoExactSecond}{0}
\newcommand{\NoOtherSecond}{0}
\newcommand{\InexactSecond}{0}
\newcommand{\StrongScored}{1,200}
\newcommand{\StrongTried}{1,824}
\newcommand{\StrongForcing}{599}
\newcommand{\StrongDecided}{25}
\newcommand{\StrongInexact}{0}
\newcommand{\ScoredAll}{909}
\newcommand{\CatalogueTable}{%
\begin{table}[ht]
\centering\footnotesize
\setlength{\tabcolsep}{3pt}
\begin{tabular}{@{}l>{\raggedright\arraybackslash}p{2.3cm}>{\raggedright\arraybackslash}p{3.3cm}rrrr@{}}
\toprule
& \multicolumn{2}{c}{The breaks the app teaches} & \multicolumn{2}{c}{His positions} & \multicolumn{2}{c}{Strong players'} \\
\cmidrule(lr){2-3}\cmidrule(lr){4-5}\cmidrule(lr){6-7}
Structure & White & Black & Break on & Scored & Drawn & Scored \\
\midrule
London pyramid & c3-c4, e3-e4 & \ldots{}c7-c5, \ldots{}e6-e5 & 638 & 427 & 438 & 178 \\
Isolated queen's pawn$^{*}$ & d4-d5 & \ldots{}b5-b4 to remove a defender & 252 & 144 & 1,171 & 408 \\
Reversed Carlsbad & f4-f5, c3-c4 & \ldots{}b5-b4, \ldots{}e6-e5 & 208 & 153 & 215 & 100 \\
Advance chain & f4-f5 & \ldots{}c7-c5, \ldots{}f7-f6 & 195 & 136 & 400 & 171 \\
Carlsbad & b4-b5, e3-e4 & \ldots{}c6-c5, \ldots{}f7-f5 & 30 & 23 & 197 & 102 \\
Pirc set-up & e4-e5, f4-f5, h4-h5 & \ldots{}c7-c5, \ldots{}e7-e5 & 18 & 11 & 229 & 96 \\
King's Indian chain & c4-c5 & \ldots{}f7-f5, \ldots{}g5-g4 & 11 & 9 & 224 & 96 \\
Hanging pawns$^{*}$ & d4-d5, c4-c5 & -- & 7 & 6 & 126 & 49 \\
\bottomrule
\end{tabular}
\caption{The app's structures with breaks, the breaks its catalogue teaches each side, and the positions with a break on: his, both halves, and the strong players' positions drawn, before the referee's quiet filter, and those scored. $^{*}$Written from the holder's side: White's column is the holder's breaks and Black's the other side's, and the app reads both with the ranks mirrored when Black holds the structure.}
\label{tab:catalogue}
\end{table}
}
\newcommand{\ArmTable}{%
\begin{table}[ht]
\centering\small
\setlength{\tabcolsep}{4pt}
\begin{tabular}{@{}llll@{}}
\toprule
& His, first half & His, second half & Strong players \\
\midrule
Positions (games) & 491 (114) & 418 (103) & 1,200 (1,200) \\
\midrule
B, the breaks & 20.5 {\scriptsize(15.9, 25.2)} & 22.4 {\scriptsize(17.2, 27.8)} & 22.8 {\scriptsize(20.5, 25.1)} \\
A, his move; S, the strong player's & 55.2 {\scriptsize(50.3, 60.1)} & 57.9 {\scriptsize(52.0, 63.5)} & 68.8 {\scriptsize(66.1, 71.4)} \\
N, Maia-2 at his band & 66.0 {\scriptsize(61.6, 70.6)} & 71.5 {\scriptsize(66.4, 76.7)} & -- \\
U, the band above & 67.6 {\scriptsize(62.9, 72.3)} & 73.0 {\scriptsize(67.9, 78.2)} & -- \\
R, a random move & 18.9 {\scriptsize(14.6, 23.2)} & 20.3 {\scriptsize(15.3, 25.6)} & 17.0 {\scriptsize(14.9, 19.2)} \\
R2, a second random move & 17.5 {\scriptsize(13.6, 21.4)} & 17.9 {\scriptsize(13.2, 22.8)} & 17.6 {\scriptsize(15.4, 19.7)} \\
\midrule
Gain: B minus A, or S (points) & $-$34.7 {\scriptsize($-$40.8, $-$28.5)} & $-$35.5 {\scriptsize($-$42.1, $-$28.7)} & $-$46.0 {\scriptsize($-$49.2, $-$42.8)} \\
B minus N (points) & $-$45.5 {\scriptsize($-$50.7, $-$39.8)} & $-$49.2 {\scriptsize($-$56.4, $-$42.3)} & -- \\
Control: A, or S, minus R (points) & 36.3 {\scriptsize(30.8, 42.0)} & 37.6 {\scriptsize(31.7, 43.4)} & 51.7 {\scriptsize(48.6, 54.7)} \\
Null: R2 minus R (points) & $-$1.4 {\scriptsize($-$5.6, 2.9)} & $-$2.4 {\scriptsize($-$6.2, 1.6)} & -- \\
\bottomrule
\end{tabular}
\caption{The share of each arm's moves that are sound, in \%, and the differences the rule reads, in points, with 95\% intervals resampling his games, or the strong players' positions. A move is sound when it gives up no more than 3 win\% against Stockfish's best.}
\label{tab:arms}
\end{table}
}
\newcommand{\MoreTable}{%
\begin{table}[ht]
\centering\small
\setlength{\tabcolsep}{4pt}
\begin{tabular}{@{}llll@{}}
\toprule
& His, first half & His, second half & Strong players \\
\midrule
Win\% given up: the breaks & 15.4 {\scriptsize(13.5, 17.3)} & 15.4 {\scriptsize(13.6, 17.2)} & 13.4 {\scriptsize(12.7, 14.1)} \\
Win\% given up: the move played & 5.6 {\scriptsize(4.8, 6.5)} & 5.5 {\scriptsize(4.5, 6.7)} & 3.1 {\scriptsize(2.8, 3.4)} \\
The breaks that are the best (\%) & 8.5 {\scriptsize(5.5, 11.8)} & 6.0 {\scriptsize(3.8, 8.5)} & 8.8 {\scriptsize(7.4, 10.4)} \\
\midrule
He, or they, played a break (\%) & 9.2 {\scriptsize(6.7, 12.0)} & 8.1 {\scriptsize(5.8, 10.7)} & 11.7 {\scriptsize(9.8, 13.5)} \\
Maia-2 at his band gives the breaks (\%) & 7.5 {\scriptsize(6.1, 9.2)} & 6.1 {\scriptsize(4.9, 7.5)} & -- \\
At the band above (\%) & 7.1 {\scriptsize(5.7, 8.9)} & 5.9 {\scriptsize(4.7, 7.3)} & -- \\
\bottomrule
\end{tabular}
\caption{The mean win\% the breaks and the move played give up against Stockfish's best (a move that scores above the best gives up none), the share of breaks that are the best, and how often a break is found: how often the move played was a break, and the probability Maia-2 gives the breaks together at his band and at the band above. 95\% intervals.}
\label{tab:more}
\end{table}
}
\newcommand{\StructureTable}{%
\begin{table}[ht]
\centering\small
\setlength{\tabcolsep}{4pt}
\begin{tabular}{@{}lrrlll@{}}
\toprule
Structure & Positions & Games & B (\%) & A (\%) & Gain (points) \\
\midrule
London pyramid & 427 & 111 & 27.6 {\scriptsize(22.2, 33.3)} & 56.4 {\scriptsize(51.0, 61.9)} & $-$28.8 {\scriptsize($-$34.9, $-$22.4)} \\
Reversed Carlsbad & 153 & 38 & 8.2 {\scriptsize(2.1, 16.5)} & 57.5 {\scriptsize(48.9, 66.0)} & $-$49.3 {\scriptsize($-$60.9, $-$37.9)} \\
Isolated queen's pawn & 144 & 48 & 22.9 {\scriptsize(16.1, 29.7)} & 54.2 {\scriptsize(46.0, 62.5)} & $-$31.3 {\scriptsize($-$40.9, $-$21.2)} \\
Advance chain & 136 & 32 & 12.5 {\scriptsize(6.2, 20.3)} & 57.4 {\scriptsize(48.7, 66.4)} & $-$44.9 {\scriptsize($-$54.5, $-$34.8)} \\
Carlsbad & 23 & 7 & \multicolumn{3}{l}{not enough data ($n=23$)} \\
Pirc set-up & 11 & 9 & \multicolumn{3}{l}{not enough data ($n=11$)} \\
King's Indian chain & 9 & 2 & \multicolumn{3}{l}{not enough data ($n=9$)} \\
Hanging pawns & 6 & 3 & \multicolumn{3}{l}{not enough data ($n=6$)} \\
\bottomrule
\end{tabular}
\caption{His positions by structure, both halves together: the share of the breaks and of his moves that are sound, and the gain, with 95\% intervals, for each structure with at least 30 scored positions.}
\label{tab:structure}
\end{table}
}
\newcommand{\SplitTable}{%
\begin{table}[ht]
\centering\small
\setlength{\tabcolsep}{4pt}
\begin{tabular}{@{}lrrlll@{}}
\toprule
His positions & Positions & Games & B (\%) & A (\%) & Gain (points) \\
\midrule
Middlegame & 819 & 209 & 20.3 {\scriptsize(17.1, 23.9)} & 56.3 {\scriptsize(52.3, 60.3)} & $-$36.0 {\scriptsize($-$40.4, $-$31.3)} \\
Endgame & 90 & 29 & 30.6 {\scriptsize(16.2, 43.5)} & 57.8 {\scriptsize(45.9, 68.3)} & $-$27.2 {\scriptsize($-$41.3, $-$14.2)} \\
\midrule
The first chance in a game & 172 & 172 & 22.1 {\scriptsize(16.9, 27.9)} & 61.0 {\scriptsize(53.5, 68.0)} & $-$39.0 {\scriptsize($-$47.4, $-$30.5)} \\
Later chances & 737 & 175 & 21.2 {\scriptsize(17.3, 25.0)} & 55.4 {\scriptsize(51.3, 59.4)} & $-$34.2 {\scriptsize($-$39.3, $-$29.2)} \\
\midrule
His move not a break, first half & 446 & 108 & 17.2 {\scriptsize(12.9, 21.9)} & 54.9 {\scriptsize(49.9, 60.2)} & $-$37.8 {\scriptsize($-$44.4, $-$31.1)} \\
His move not a break, second half & 384 & 96 & 20.1 {\scriptsize(14.9, 25.5)} & 57.8 {\scriptsize(52.1, 63.4)} & $-$37.8 {\scriptsize($-$44.6, $-$31.0)} \\
\bottomrule
\end{tabular}
\caption{His positions by the divider's phase and by whether the position is the first in its game with a break on, both halves together, and his positions where his move was not a break, by half: the share of the breaks and of his moves that are sound, and the gain, with 95\% intervals.}
\label{tab:splits}
\end{table}
}
\newcommand{\Verdict}{the break misleads}
\newcommand{\BFirst}{20.5}
\newcommand{\AFirst}{55.2}
\newcommand{\NFirst}{66.0}
\newcommand{\UFirst}{67.6}
\newcommand{\RFirst}{18.9}
\newcommand{\RTwoFirst}{17.5}
\newcommand{\GainFirst}{$-$34.7}
\newcommand{\GainFirstCi}{$-$40.8 to $-$28.5}
\newcommand{\BNFirst}{$-$45.5}
\newcommand{\BNCiFirst}{$-$50.7 to $-$39.8}
\newcommand{\ControlFirst}{36.3}
\newcommand{\ControlFirstCi}{30.8 to 42.0}
\newcommand{\NullFirst}{$-$1.4}
\newcommand{\NullFirstCi}{$-$5.6 to 2.9}
\newcommand{\LostBFirst}{15.4}
\newcommand{\LostAFirst}{5.6}
\newcommand{\BestBFirst}{8.5}
\newcommand{\PlayedFirst}{9.2}
\newcommand{\MassHisFirst}{7.5}
\newcommand{\MassPlusFirst}{7.1}
\newcommand{\BSecond}{22.4}
\newcommand{\ASecond}{57.9}
\newcommand{\NSecond}{71.5}
\newcommand{\USecond}{73.0}
\newcommand{\RSecond}{20.3}
\newcommand{\RTwoSecond}{17.9}
\newcommand{\GainSecond}{$-$35.5}
\newcommand{\GainSecondCi}{$-$42.1 to $-$28.7}
\newcommand{\BNSecond}{$-$49.2}
\newcommand{\BNCiSecond}{$-$56.4 to $-$42.3}
\newcommand{\ControlSecond}{37.6}
\newcommand{\ControlSecondCi}{31.7 to 43.4}
\newcommand{\NullSecond}{$-$2.4}
\newcommand{\NullSecondCi}{$-$6.2 to 1.6}
\newcommand{\LostBSecond}{15.4}
\newcommand{\LostASecond}{5.5}
\newcommand{\BestBSecond}{6.0}
\newcommand{\PlayedSecond}{8.1}
\newcommand{\MassHisSecond}{6.1}
\newcommand{\MassPlusSecond}{5.9}
\newcommand{\StrongB}{22.8}
\newcommand{\StrongS}{68.8}
\newcommand{\StrongR}{17.0}
\newcommand{\StrongRTwo}{17.6}
\newcommand{\StrongGain}{$-$46.0}
\newcommand{\StrongGainCi}{$-$49.2 to $-$42.8}
\newcommand{\StrongControl}{51.7}
\newcommand{\StrongControlCi}{48.6 to 54.7}
\newcommand{\StrongLostB}{13.4}
\newcommand{\StrongLostS}{3.1}
\newcommand{\StrongBestB}{8.8}
\newcommand{\StrongPlayed}{11.7}
\newcommand{\StructLondonPyramidB}{27.6}
\newcommand{\StructLondonPyramidA}{56.4}
\newcommand{\StructLondonPyramidGain}{$-$28.8}
\newcommand{\StructLondonPyramidN}{427}
\newcommand{\StructCarlsbadReversedB}{8.2}
\newcommand{\StructCarlsbadReversedA}{57.5}
\newcommand{\StructCarlsbadReversedGain}{$-$49.3}
\newcommand{\StructCarlsbadReversedN}{153}
\newcommand{\StructIqpB}{22.9}
\newcommand{\StructIqpA}{54.2}
\newcommand{\StructIqpGain}{$-$31.3}
\newcommand{\StructIqpN}{144}
\newcommand{\StructAdvanceChainB}{12.5}
\newcommand{\StructAdvanceChainA}{57.4}
\newcommand{\StructAdvanceChainGain}{$-$44.9}
\newcommand{\StructAdvanceChainN}{136}
\newcommand{\MidB}{20.3}
\newcommand{\MidA}{56.3}
\newcommand{\MidGain}{$-$36.0}
\newcommand{\EndB}{30.6}
\newcommand{\EndA}{57.8}
\newcommand{\EndGain}{$-$27.2}
\newcommand{\EndN}{90}
\newcommand{\FirstB}{22.1}
\newcommand{\FirstA}{61.0}
\newcommand{\FirstGain}{$-$39.0}
\newcommand{\LaterB}{21.2}
\newcommand{\LaterA}{55.4}
\newcommand{\LaterGain}{$-$34.2}
\newcommand{\NotPlayedFirst}{446}
\newcommand{\NotPlayedGainFirst}{$-$37.8}
\newcommand{\NotPlayedBFirst}{17.2}
\newcommand{\NotPlayedSecond}{384}
\newcommand{\NotPlayedGainSecond}{$-$37.8}
\newcommand{\NotPlayedBSecond}{20.1}
\newcommand{\SomeFirst}{26.3}
\newcommand{\SomeMinusFirst}{$-$28.9}
\newcommand{\SomeMinusFirstCi}{$-$35.5 to $-$21.9}
\newcommand{\SomeSecond}{28.0}
\newcommand{\SomeMinusSecond}{$-$29.9}
\newcommand{\SomeMinusSecondCi}{$-$37.0 to $-$22.4}
\newcommand{\BRFirst}{1.5}
\newcommand{\BRFirstCi}{$-$3.5 to 6.6}
\newcommand{\BRSecond}{2.0}
\newcommand{\BRSecondCi}{$-$3.8 to 7.8}
\newcommand{\StrongBR}{5.8}
\newcommand{\StrongBRCi}{2.9 to 8.6}
\newcommand{\StrongSome}{26.3}
\newcommand{\StrongSomeMinus}{$-$42.4}
\newcommand{\StrongSomeMinusCi}{$-$45.8 to $-$39.0}
\newcommand{\PlayedN}{79}
\newcommand{\PlayedSound}{58.2}
\newcommand{\PlayedSoundCi}{47.4 to 69.4}
\newcommand{\NotPlayedB}{18.5}
\newcommand{\NotPlayedBCi}{15.2 to 22.0}
\newcommand{\StrongPlayedN}{140}
\newcommand{\StrongPlayedSound}{70.0}
\newcommand{\StrongNotPlayedB}{17.8}
\newcommand{\BestN}{94}
\newcommand{\BestShare}{10.3}
\newcommand{\BestPlayed}{25.5}
\newcommand{\BestPlayedCi}{17.1 to 33.7}
\newcommand{\NotBestPlayed}{6.7}
\newcommand{\BestMass}{19.5}
\newcommand{\NotBestMass}{5.4}
\newcommand{\BestMassPlus}{20.1}
\newcommand{\NotBestMassPlus}{5.0}
\newcommand{\StrongBestN}{132}
\newcommand{\StrongBestPlayed}{42.4}
\newcommand{\StrongNotBestPlayed}{7.9}
\newcommand{\MomentTable}{%
\begin{table}[ht]
\centering\footnotesize
\setlength{\tabcolsep}{3pt}
\begin{tabular}{@{}lrllll@{}}
\toprule
Positions & $n$ & \begin{tabular}[b]{@{}l@{}}Played \\ a break\end{tabular} & \begin{tabular}[b]{@{}l@{}}Played \\ the best\end{tabular} & \begin{tabular}[b]{@{}l@{}}Maia-2 gives the \\ breaks, his band\end{tabular} & \begin{tabular}[b]{@{}l@{}}The band \\ above\end{tabular} \\
\midrule
His, a break is the best & 94 & 25.5 {\scriptsize(17.1, 33.7)} & 23.4 {\scriptsize(14.8, 32.0)} & 19.5 {\scriptsize(15.5, 24.7)} & 20.1 {\scriptsize(15.9, 25.5)} \\
His, no break is the best & 815 & 6.7 {\scriptsize(5.2, 8.4)} & 22.5 {\scriptsize(19.6, 25.4)} & 5.4 {\scriptsize(4.6, 6.2)} & 5.0 {\scriptsize(4.2, 5.8)} \\
Strong, a break is the best & 132 & 42.4 {\scriptsize(34.1, 50.8)} & 40.2 {\scriptsize(31.8, 48.5)} & -- & -- \\
Strong, no break is the best & 1,068 & 7.9 {\scriptsize(6.3, 9.5)} & 36.0 {\scriptsize(33.1, 39.1)} & -- & -- \\
\bottomrule
\end{tabular}
\caption{Added after the run. His positions, both halves together, and the strong players', where a break is Stockfish's best and where none is: how often the move played was a break and was the best, and the probability Maia-2 gives the breaks together. In \%, with 95\% intervals.}
\label{tab:moment}
\end{table}
}
\newcommand{\ChoiceTable}{%
\begin{table}[ht]
\centering\footnotesize
\setlength{\tabcolsep}{3pt}
\begin{tabular}{@{}lrllll@{}}
\toprule
Positions & $n$ & \begin{tabular}[b]{@{}l@{}}Some break \\ sound\end{tabular} & \begin{tabular}[b]{@{}l@{}}The move \\ played, sound\end{tabular} & \begin{tabular}[b]{@{}l@{}}Difference \\ (points)\end{tabular} & \begin{tabular}[b]{@{}l@{}}B minus R \\ (points)\end{tabular} \\
\midrule
His, first half & 491 & 26.3 {\scriptsize(21.0, 31.9)} & 55.2 {\scriptsize(50.3, 60.1)} & $-$28.9 {\scriptsize($-$35.5, $-$21.9)} & 1.5 {\scriptsize($-$3.5, 6.6)} \\
His, second half & 418 & 28.0 {\scriptsize(21.4, 34.9)} & 57.9 {\scriptsize(52.1, 64.0)} & $-$29.9 {\scriptsize($-$37.0, $-$22.4)} & 2.0 {\scriptsize($-$3.8, 7.8)} \\
Strong players & 1,200 & 26.3 {\scriptsize(23.8, 28.9)} & 68.8 {\scriptsize(66.1, 71.4)} & $-$42.4 {\scriptsize($-$45.8, $-$39.0)} & 5.8 {\scriptsize(2.9, 8.6)} \\
\midrule
& & \begin{tabular}[b]{@{}l@{}}The break \\ played, sound\end{tabular} & \begin{tabular}[b]{@{}l@{}}The breaks, \\ sound\end{tabular} & \begin{tabular}[b]{@{}l@{}}The move \\ played, sound\end{tabular} & \\
\midrule
His, played one & 79 & 58.2 {\scriptsize(47.4, 69.4)} & 51.3 {\scriptsize(41.9, 61.3)} & -- & -- \\
His, did not & 830 & -- & 18.5 {\scriptsize(15.2, 22.0)} & 56.3 {\scriptsize(52.3, 60.2)} & -- \\
Strong, played one & 140 & 70.0 {\scriptsize(62.1, 77.1)} & 60.0 {\scriptsize(53.2, 66.8)} & -- & -- \\
Strong, did not & 1,060 & -- & 17.8 {\scriptsize(15.7, 20.0)} & 68.6 {\scriptsize(65.8, 71.3)} & -- \\
\bottomrule
\end{tabular}
\caption{Added after the run. Above, the share of positions where at least one break is sound, against the share where the move played is, and the breaks' soundness minus a random move's, paired by position. Below, his positions, both halves together, and the strong players', where the move played was a break or was not: the share of the breaks played that are sound, of all the breaks, and of the moves played. In \%, with 95\% intervals.}
\label{tab:choice}
\end{table}
}
\newcommand{\LeverTable}{%
\begin{table}[ht]
\centering\footnotesize
\setlength{\tabcolsep}{3pt}
\begin{tabular}{@{}>{\raggedright\arraybackslash}p{2.4cm}>{\raggedright\arraybackslash}p{2.3cm}rlrrlr@{}}
\toprule
& & \multicolumn{3}{c}{His positions} & \multicolumn{3}{c}{Strong players'} \\
\cmidrule(lr){3-5}\cmidrule(lr){6-8}
Structure & Break & $n$ & The break & His move & $n$ & The break & Their move \\
\midrule
London pyramid & c3-c4 & 373 & 29.0 {\scriptsize(22.6, 35.4)} & 56.6 & 63 & 25.4 {\scriptsize(15.9, 36.5)} & 74.6 \\
London pyramid & e3-e4 & 355 & 23.4 {\scriptsize(17.7, 29.5)} & 56.3 & 80 & 17.5 {\scriptsize(10.0, 26.3)} & 75.0 \\
Reversed Carlsbad & \ldots{}e6-e5 & 151 & 7.9 {\scriptsize(2.0, 15.6)} & 58.3 & 48 & 12.5 {\scriptsize(4.2, 22.9)} & 66.7 \\
Advance chain & \ldots{}f7-f6 & 118 & 9.3 {\scriptsize(3.3, 18.1)} & 58.5 & 132 & 10.6 {\scriptsize(6.1, 15.9)} & 62.1 \\
Isolated queen's pawn & d5-d4 & 76 & 25.0 {\scriptsize(16.0, 35.1)} & 50.0 & 159 & 22.6 {\scriptsize(16.4, 28.9)} & 64.8 \\
Isolated queen's pawn & d4-d5 & 44 & 11.4 {\scriptsize(2.3, 21.1)} & 50.0 & 143 & 22.4 {\scriptsize(16.1, 29.4)} & 68.5 \\
Carlsbad & e3-e4 & 22 & -- & -- & 55 & 12.7 {\scriptsize(5.5, 21.8)} & 70.9 \\
London pyramid & \ldots{}c7-c5 & 22 & -- & -- & 39 & 28.2 {\scriptsize(15.4, 43.6)} & 71.8 \\
Advance chain & \ldots{}c7-c5 & 18 & -- & -- & 37 & 32.4 {\scriptsize(18.9, 48.6)} & 70.3 \\
Isolated queen's pawn & \ldots{}b4-b5 to remove a defender & 16 & -- & -- & 42 & 31.0 {\scriptsize(16.7, 45.2)} & 73.8 \\
London pyramid & \ldots{}e6-e5 & 12 & -- & -- & 66 & 15.2 {\scriptsize(7.6, 24.2)} & 72.7 \\
King's Indian chain & \ldots{}f7-f5 & 9 & -- & -- & 51 & 43.1 {\scriptsize(29.4, 56.9)} & 70.6 \\
Isolated queen's pawn & \ldots{}b5-b4 to remove a defender & 8 & -- & -- & 64 & 25.0 {\scriptsize(15.6, 35.9)} & 60.9 \\
Pirc set-up & \ldots{}e7-e5 & 8 & -- & -- & 32 & 34.4 {\scriptsize(18.8, 50.0)} & 75.0 \\
Pirc set-up & e4-e5 & 3 & -- & -- & 47 & 23.4 {\scriptsize(12.8, 36.2)} & 78.7 \\
Reversed Carlsbad & c3-c4 & 2 & -- & -- & 50 & 14.0 {\scriptsize(6.0, 24.0)} & 74.0 \\
Pirc set-up & \ldots{}c7-c5 & 2 & -- & -- & 36 & 38.9 {\scriptsize(25.0, 55.6)} & 77.8 \\
Carlsbad & \ldots{}c6-c5 & 1 & -- & -- & 47 & 21.3 {\scriptsize(10.6, 34.0)} & 76.6 \\
King's Indian chain & c4-c5 & 0 & -- & -- & 36 & 33.3 {\scriptsize(19.4, 50.0)} & 55.6 \\
\bottomrule
\end{tabular}
\caption{Added after the run. Each break the catalogue teaches, apart, where it is on in at least 30 of his positions (both halves) or of the strong players': the positions, the share of the break's moves that are sound, with its 95\% interval, and the share of the moves played there that are, in \%. Each break is the catalogue's text as the app reads it for the side to move: for a structure written from the holder's side, with the ranks mirrored when Black holds the structure.}
\label{tab:levers}
\end{table}
}
\newcommand{\LevOneStructure}{the London pyramid}
\newcommand{\LevOneBreak}{c3-c4}
\newcommand{\LevOneN}{373}
\newcommand{\LevOneL}{29.0}
\newcommand{\LevOneOwn}{56.6}
\newcommand{\LevTwoStructure}{the London pyramid}
\newcommand{\LevTwoBreak}{e3-e4}
\newcommand{\LevTwoN}{355}
\newcommand{\LevTwoL}{23.4}
\newcommand{\LevTwoOwn}{56.3}
\newcommand{\LevThreeStructure}{the reversed Carlsbad}
\newcommand{\LevThreeBreak}{\ldots{}e6-e5}
\newcommand{\LevThreeN}{151}
\newcommand{\LevThreeL}{7.9}
\newcommand{\LevThreeOwn}{58.3}
\newcommand{\LevFourStructure}{the advance chain}
\newcommand{\LevFourBreak}{\ldots{}f7-f6}
\newcommand{\LevFourN}{118}
\newcommand{\LevFourL}{9.3}
\newcommand{\LevFourOwn}{58.5}
\newcommand{\StrongLevOneStructure}{the isolated queen's pawn}
\newcommand{\StrongLevOneBreak}{d5-d4}
\newcommand{\StrongLevOneN}{159}
\newcommand{\StrongLevOneL}{22.6}
\newcommand{\StrongLevOneOwn}{64.8}
\newcommand{\StrongLevTwoStructure}{the isolated queen's pawn}
\newcommand{\StrongLevTwoBreak}{d4-d5}
\newcommand{\StrongLevTwoN}{143}
\newcommand{\StrongLevTwoL}{22.4}
\newcommand{\StrongLevTwoOwn}{68.5}
\newcommand{\StrongLevThreeStructure}{the advance chain}
\newcommand{\StrongLevThreeBreak}{\ldots{}f7-f6}
\newcommand{\StrongLevThreeN}{132}
\newcommand{\StrongLevThreeL}{10.6}
\newcommand{\StrongLevThreeOwn}{62.1}
\newcommand{\StrongLevFourStructure}{the London pyramid}
\newcommand{\StrongLevFourBreak}{e3-e4}
\newcommand{\StrongLevFourN}{80}
\newcommand{\StrongLevFourL}{17.5}
\newcommand{\StrongLevFourOwn}{75.0}
\newcommand{\StrongStructureTable}{%
\begin{table}[ht]
\centering\small
\setlength{\tabcolsep}{4pt}
\begin{tabular}{@{}lrlll@{}}
\toprule
Structure & Positions & B (\%) & S (\%) & B minus S (points) \\
\midrule
Isolated queen's pawn & 408 & 23.8 {\scriptsize(19.9, 27.9)} & 66.4 {\scriptsize(62.0, 70.8)} & $-$42.6 {\scriptsize($-$48.0, $-$37.0)} \\
London pyramid & 178 & 19.7 {\scriptsize(14.9, 25.0)} & 72.5 {\scriptsize(65.7, 79.2)} & $-$52.8 {\scriptsize($-$60.1, $-$44.7)} \\
Advance chain & 171 & 15.2 {\scriptsize(10.5, 20.5)} & 65.5 {\scriptsize(58.5, 72.5)} & $-$50.3 {\scriptsize($-$58.8, $-$42.1)} \\
Carlsbad & 102 & 16.7 {\scriptsize(9.8, 23.5)} & 73.5 {\scriptsize(64.7, 82.4)} & $-$56.9 {\scriptsize($-$66.7, $-$46.6)} \\
Reversed Carlsbad & 100 & 15.0 {\scriptsize(9.0, 22.0)} & 71.0 {\scriptsize(62.0, 80.0)} & $-$56.0 {\scriptsize($-$68.0, $-$43.5)} \\
King's Indian chain & 96 & 38.0 {\scriptsize(28.1, 47.9)} & 63.5 {\scriptsize(54.2, 72.9)} & $-$25.5 {\scriptsize($-$37.5, $-$13.5)} \\
Pirc set-up & 96 & 33.9 {\scriptsize(25.0, 42.7)} & 77.1 {\scriptsize(68.8, 85.4)} & $-$43.2 {\scriptsize($-$55.2, $-$31.3)} \\
Hanging pawns & 49 & 28.6 {\scriptsize(21.4, 35.7)} & 65.3 {\scriptsize(51.0, 79.6)} & $-$36.7 {\scriptsize($-$51.0, $-$22.4)} \\
\bottomrule
\end{tabular}
\caption{Added after the run. The strong players' positions by structure: the share of the breaks and of their moves that are sound, and the difference, with 95\% intervals, for each structure with at least 30 scored positions.}
\label{tab:strongstructure}
\end{table}
}
\newcommand{\EaHalf}{second}
\newcommand{\EbHalf}{second}
\newcommand{\LondonWhite}{601}
\newcommand{\EaFen}{\fen{r2q1rk1/ppp1ppbp/2np1np1/8/2PPP1bP/2N1BN2/PP3PP1/R2QKB1R b KQ - 2 8}}
\newcommand{\EaMover}{Black}
\newcommand{\EaPhase}{middlegame}
\newcommand{\EaStructure}{the Pirc set-up}
\newcommand{\EaBreaks}{8\ldots{}e5 (sound, 55\%)}
\newcommand{\EaBreakLabel}{The break}
\newcommand{\EaBest}{8\ldots{}e5}
\newcommand{\EaWpBest}{55}
\newcommand{\EaPlayed}{8\ldots{}e6}
\newcommand{\EaWpPlayed}{49}
\newcommand{\EaPlayedSound}{not sound}
\newcommand{\EaLost}{$-$6}
\newcommand{\EaBoard}{%
\chessboard[setfen={r2q1rk1/ppp1ppbp/2np1np1/8/2PPP1bP/2N1BN2/PP3PP1/R2QKB1R b KQ - 2 8}, boardfontsize=14pt, pgfstyle=straightmove, arrow=to, linewidth=0.07em, color=blue!70!black, markmoves={e7-e6}, pgfstyle=straightmove, arrow=to, linewidth=0.07em, color=green!50!black, markmoves={e7-e5}]
}
\newcommand{\EbFen}{\fen{rn1q2k1/ppp3p1/3b1p1p/3p1R1Q/3Pr2B/2P1P3/PP4PP/RN4K1 w - - 0 16}}
\newcommand{\EbMover}{White}
\newcommand{\EbPhase}{middlegame}
\newcommand{\EbStructure}{the London pyramid}
\newcommand{\EbBreaks}{16.~c4 (not sound, 54\%)}
\newcommand{\EbBreakLabel}{The break}
\newcommand{\EbBest}{16.~\symknight{}d2}
\newcommand{\EbWpBest}{61}
\newcommand{\EbPlayed}{16.~\symknight{}d2}
\newcommand{\EbWpPlayed}{61}
\newcommand{\EbPlayedSound}{sound}
\newcommand{\EbLost}{0}
\newcommand{\EbBoard}{%
\chessboard[setfen={rn1q2k1/ppp3p1/3b1p1p/3p1R1Q/3Pr2B/2P1P3/PP4PP/RN4K1 w - - 0 16}, boardfontsize=14pt, pgfstyle=straightmove, arrow=to, linewidth=0.07em, color=blue!70!black, markmoves={b1-d2}, pgfstyle=straightmove, arrow=to, linewidth=0.07em, color=red!75!black, markmoves={c3-c4}]
}


\title{Do the breaks the app teaches for each pawn structure hold at his level?\\[0.3em]
\large Each structure's breaks against his own move in \ScoredAll{} of his quiet positions,
refereed by Stockfish, with strong players' positions beside them}
\author{Chess Coach research lab}
\date{}

\begin{document}
\maketitle

\begin{abstract}
When the app names the pawn structure in one of his positions and a break its catalogue teaches
for that structure can be played, is the break a sound move there, more often than the move he
played? In \ScoredAll{} quiet middlegame and endgame positions from his 10-minute games, Stockfish
at \GradeNodes{} nodes calls a move sound when it gives up no more than \Window{} win\% against
the best. By the rule fixed before any move was scored, the break misleads: the breaks are sound
\BFirst\% and \BSecond\% of the time on the two halves of his games, his own moves \AFirst\% and
\ASecond\%, a gain of \GainFirst{} and \GainSecond{} points with both intervals below 0, and both
controls pass. A break that can be played is sound about as often as a random move that is not a
break, and gives up \LostBFirst{} and \LostBSecond{} win\% on average against his moves'
\LostAFirst{} and \LostASecond. It is so in every structure with enough positions, and in
\StrongScored{} positions of players rated \Floor{} or more, whose breaks are sound \StrongB\% of
the time against their own moves' \StrongS\%: a break is rarely right merely because it can be
played. Added after the run: the breaks he played were sound \PlayedSound\% of the time, the ones
he passed over \NotPlayedB\%, and where a break is Stockfish's best, in \BestShare\% of these
positions, he played one \BestPlayed\% of the time against \NotBestPlayed\% elsewhere. The app
should not offer a structure's break as the plan whenever it is on; what it can teach is the
moment, with the engine's word on each position.
\end{abstract}

\section{The question}

A coach explains a middlegame by its pawn structure. In the Carlsbad the plan is the minority
attack and its break is b4-b5; against the London pyramid, Black strikes with \ldots c5. The app
does the same. It names the structure on the board, teaches each structure's plans and breaks
from its catalogue, and its structure index counts the games in which he played the structure's
break. This study asks whether that teaching holds in his games: when a break his structure
teaches can be played in one of his quiet middlegame or endgame positions, is it a sound move
there, more often than the move he played?

If the break holds, the review and the drills can teach the structure's break as the plan when it
is on, and the count of breaks played measures something worth raising. If it misleads, playing
the break whenever it can be played is worse than his own choice, and the app has to teach when a
break is right, not only which break.

This study follows the one on positions like his. There, strong players' moves from positions
like his carried over weakly, and best where the app names the pawn structure, which suggested
testing the plan a structure calls for in place of the next move from a similar game.

\section{How it was measured}

\paragraph{His positions.} His \Games{} 10-minute games in the app's archive at commit
\texttt{\AppCommit}, in order of their end time: the first \GamesFirst{} (\FirstFrom{} to
\FirstTo) and the last \GamesSecond{} (\SecondFrom{} to \SecondTo) are the two halves. Nothing is
fitted, so the halves are two replications. The positions are the ones the study on positions like
his drew: before one of his moves, past the opening by the app's divider, not in check, with at
least \MinLegal{} legal moves and \MinMen{} men or more, his win\% before the move between
\WinLow{} and \WinHigh{} by the archive's evaluation, and his clock at least \ClockFloor{} seconds,
\Candidates{} positions in all. The app names a structure in \Named{} of them (\NamedShare\%). A
position enters this study when a break is on: at least one of his legal moves is a break of that
structure for his side. That leaves \BreakOn{} positions from \BreakGames{} games, \BreakOnFirst{}
in the first half and \BreakOnSecond{} in the second. In \NotOn{} positions with a named
structure no break was on, and in \NoLever{} the catalogue gives his side no break.

\paragraph{The breaks.} The app's catalogue writes each structure's breaks for each side as pawn
moves, such as c3-c4 (Table~\ref{tab:catalogue}). For three structures, the isolated queen's
pawn, hanging pawns and the queenside majority, it writes them from the holder's side, the side
whose pawns make the structure, and the app reads them with the ranks mirrored when Black holds
it. A move is a break when the app's own reading counts it, played after the game's moves so far:
a pawn move of his to the last square a break names, a capture landing there included. The
structure is the one the app names on the board of the position.

\paragraph{Strong players' positions.} The study on positions like his streamed the first
\StreamCap{} rated games of Lichess's \MonthName{} database (CC0) between two players rated
\Floor{} or more. Read in file order, each game offers the positions before a move by either
side that are past the opening, not in check, with \MinMen{} men or more, at least \MinLegal{}
legal moves, the mover's clock at least \ClockFloor{} seconds and a structure the app names. A
seeded hash of the game shuffles them, and the first of up to \Tries{} with a break on for the
mover is the game's one position. The first \StrongGamesRead{} games gave \StrongDrawn{}
positions; \StrongNoStructure{} games had no position with a named structure, and in
\StrongNotOn{} none of the tries had a break on. Most are blitz: \Blitz{} blitz, \Rapid{} rapid
and \Classical{} classical. They are scored in order until \StrongTarget{} pass the referee's
quiet filter.

\paragraph{The moves compared.} In each of his positions, \textbf{B} is every legal move of his
that is a break, and B's soundness is the share of them that are sound. \textbf{A} is his move.
\textbf{N} is Maia-2's most likely move at his band and \textbf{U} at the band above: his band is
his rating at the game with the app's offset of \Offset{} points, the band above the same of his
rating plus \Above, each conditioned on the opponent's band. \textbf{R} and \textbf{R2} are two
legal moves that are not breaks, each drawn uniformly at random with a seed, independently of each
other. In the strong players' positions: B, \textbf{S}, the strong player's move, R and R2.

\paragraph{The referee.} Stockfish through the app's engine service at the grading budget of
\GradeNodes{} nodes, with its hash cleared before each position and no tablebase, gives the best
move and scores every move compared. A move is \emph{sound} when it gives up no more than
\Window{} win\% against the best, by the app's own test for an alternative move. A position is
\emph{quiet} when the best move is not a capture, a check or a promotion and the win\% it gives
the side to move is between \WinLow{} and \WinHigh; only quiet positions are scored. A score
that is not exact counts as missing, and a position with one leaves the analysis.

\paragraph{The measures.} The \emph{gain} is B's soundness minus A's, averaged over positions,
paired by position; where his move was a break, B and A share it. \Draws{} bootstrap draws with a
fixed seed resample his games within each half, or the strong players' positions, and intervals
are the 95\% percentile intervals.

\paragraph{The rule, fixed before any move was scored.} Two controls come first, and if either
fails the study gives no verdict.
\begin{enumerate}[nosep]
\item \textbf{The referee tells moves apart here.} His own move's soundness minus R's must have
  its interval above 0 on each half, and the strong player's move's minus R's on the strong
  players' positions.
\item \textbf{The null.} R2 against R, in place of B against A, must read neither as the break
  holding nor as it misleading.
\end{enumerate}
Then, on his positions, \emph{the break holds} when the gain's interval lies above 0 on both
halves, \emph{the break misleads} when it lies below 0 on both, and it is \emph{no better than his
own move} otherwise. A half with fewer than \MinPositions{} scored positions is \emph{not enough
data}, whatever the numbers say. The rule and the scripts went to the lab before any move was
scored; the positions with a break on had been drawn and counted, with no move scored.

\section{Coverage}

\CatalogueTable

A break is on in \OnShare\% of his positions with a named structure (Table~\ref{tab:catalogue}).
\LondonShare\% of them are the London pyramid, \LondonWhite{} of its \LondonOn{} with him as White,
and nearly all of his reversed Carlsbad and advance chain positions have him as Black. The strong
players' positions are most often the isolated queen's pawn, \StrongIqpShare\% of those drawn.

The referee keeps the quiet ones: \ScoredFirst{} positions from \ScoredGamesFirst{} games in the
first half and \ScoredSecond{} from \ScoredGamesSecond{} games in the second. In \ForcingFirst{}
and \ForcingSecond{} positions the best move is a capture, a check or a promotion, and in
\DecidedFirst{} and \DecidedSecond{} the referee reads the game as decided. \StrongTried{} strong
players' positions were tried before \StrongScored{} were quiet. Every score was exact.

\section{The controls}

\paragraph{The referee tells moves apart.} His own move is sound \ControlFirst{} and
\ControlSecond{} points more often than a random move that is not a break (\ControlFirstCi;
\ControlSecondCi), and the strong player's move \StrongControl{} points more often
(\StrongControlCi). All three intervals lie above 0.

\paragraph{The null passes.} R2 against R reads \NullFirst{} points (\NullFirstCi) in the first
half and \NullSecond{} points (\NullSecondCi) in the second; both intervals hold 0.

\section{The verdict: the break misleads}

\ArmTable

The breaks are sound \BFirst\% and \BSecond\% of the time on the two halves, his own moves
\AFirst\% and \ASecond\% (Table~\ref{tab:arms}). The gain is \GainFirst{} points (\GainFirstCi)
and \GainSecond{} points (\GainSecondCi), below 0 on both halves, so by the rule the break
misleads: playing the structure's break whenever it can be played would be worse than his own
choice.

A break that can be played is sound about as often as a random move that is not a break, at
\RFirst\% and \RSecond\%; paired by position, the breaks are \BRFirst{} points (\BRFirstCi) and
\BRSecond{} points (\BRSecondCi) above it (Table~\ref{tab:choice}). Being one of the structure's
breaks adds nothing to a move's chance of being sound when all that is known is that it can be
played. Maia-2's most likely move at his band is sound \NFirst\% and \NSecond\% of the time, and
at the band above \UFirst\% and \USecond\%, more often than his own.

\section{What the positions describe}

\MoreTable

\paragraph{How much a break gives up.} On average the breaks give up \LostBFirst{} and
\LostBSecond{} win\% against the best, about three times what his moves give up, \LostAFirst{}
and \LostASecond{} (Table~\ref{tab:more}). \BestBFirst\% and \BestBSecond\% of the breaks are
the referee's best.

\paragraph{Finding it.} He played a break in \PlayedFirst\% and \PlayedSecond\% of these
positions, and Maia-2 at his band gives the breaks together \MassHisFirst\% and \MassHisSecond\%
of its probability, at the band above \MassPlusFirst\% and \MassPlusSecond\%. Strong players
played a break in \StrongPlayed\% of theirs.

\paragraph{Strong players.} In \StrongScored{} positions of players rated \Floor{} or more, mostly
blitz, the breaks are sound \StrongB\% of the time against the strong player's move's \StrongS\%,
\StrongGain{} points (\StrongGainCi), and a random move that is not a break \StrongR\%; paired by
position, their breaks are \StrongBR{} points above it (\StrongBRCi). The break that can be played
is rarely the move at their level either, and trails their own move by more than his trails his,
so what misleads is the reading of a break as the plan whenever it is on, not his level.

\section{By structure, phase and chance}

\StructureTable

\paragraph{Every structure.} The London pyramid's breaks are sound \StructLondonPyramidB\% of the
time against his move's \StructLondonPyramidA\%, the isolated queen's pawn's \StructIqpB\% against
\StructIqpA\%, the reversed Carlsbad's \StructCarlsbadReversedB\% against
\StructCarlsbadReversedA\%, and the advance chain's \StructAdvanceChainB\% against
\StructAdvanceChainA\% (Table~\ref{tab:structure}). In each structure with enough positions the
gain's interval lies below 0: it is not only the London.

\SplitTable

\paragraph{Phase, the first chance, and the breaks he passed over.} In the middlegame the breaks
are sound \MidB\% of the time against his move's \MidA\%, and in the endgame, in only \EndN{}
positions, \EndB\% against \EndA\% (Table~\ref{tab:splits}). The first position in a game with a
break on is no better for the break than the later ones, \FirstB\% against \LaterB\%. Where his move was not a break, the breaks he
passed over are sound \NotPlayedBFirst\% and \NotPlayedBSecond\% of the time on the two halves.

\section{Added after the run}

These descriptions were added once his halves had been scored and looked at; none was named before
the run, and they decide nothing.

\MomentTable

\paragraph{The moment.} A break is Stockfish's best in \BestN{} of his positions, \BestShare\% of
them. There he played a break \BestPlayed\% of the time (\BestPlayedCi), against \NotBestPlayed\%
where no break is the best, and Maia-2 at his band gives the breaks \BestMass\% against
\NotBestMass\% (Table~\ref{tab:moment}). He and his band read when a break is due, and both miss it
most of the time. Strong players played a break \StrongBestPlayed\% of the time where one is the
best, in \StrongBestN{} positions, and \StrongNotBestPlayed\% elsewhere.

\ChoiceTable

\paragraph{His own choice.} Where he played a break, in \PlayedN{} positions, it was sound
\PlayedSound\% of the time (\PlayedSoundCi), about as often as his moves are; where he did not,
the breaks he passed over were sound \NotPlayedB\% (\NotPlayedBCi). He breaks more selectively
than the catalogue's reading of when a break is on (Table~\ref{tab:choice}). Strong players do the same: the breaks they
played, in \StrongPlayedN{} positions, were sound \StrongPlayedSound\% of the time, the ones they
passed over \StrongNotPlayedB\%. At least one break is sound in \SomeFirst\% and \SomeSecond\% of
his positions, against \AFirst\% and \ASecond\% for his move.

\LeverTable

\paragraph{Each break apart.} His most common breaks are \LevOneStructure's \LevOneBreak, sound
\LevOneL\% of the time in \LevOneN{} positions against his move's \LevOneOwn\%,
\LevTwoStructure's \LevTwoBreak{} (\LevTwoL\% against \LevTwoOwn\%), \LevThreeStructure's
\LevThreeBreak{} (\LevThreeL\% against \LevThreeOwn\%) and \LevFourStructure's \LevFourBreak{}
(\LevFourL\% against \LevFourOwn\%) (Table~\ref{tab:levers}). No break with enough positions, his
or the strong players', is sound as often as the move played. The strong players' most common is
\StrongLevOneStructure's \StrongLevOneBreak, sound \StrongLevOneL\% of the time against their
move's \StrongLevOneOwn\%.

\StrongStructureTable

\section{Two positions}

Two of his positions show what a break that can be played looks like (Figure~\ref{fig:examples}):
the first of the \EaHalf{} half where a break is the best move and he played neither a break nor
another sound move, and the first London pyramid position of the \EbHalf{} half where every break
is unsound and his own move is sound. They illustrate; they test nothing.

\begin{figure}[ht]
\centering
\begin{subfigure}[t]{0.48\textwidth}
\centering
\EaBoard\\[0.3em]
{\footnotesize FEN \EaFen}
\caption{He is \EaMover, in the \EaPhase, in \EaStructure. The best move is \EaBest, at
\EaWpBest\% to win. \EaBreakLabel: \EaBreaks. He played \EaPlayed{} (blue), \EaPlayedSound{} at
\EaWpPlayed\%.}
\label{fig:due}
\end{subfigure}\hfill
\begin{subfigure}[t]{0.48\textwidth}
\centering
\EbBoard\\[0.3em]
{\footnotesize FEN \EbFen}
\caption{He is \EbMover, in the \EbPhase, in \EbStructure. The best move is \EbBest, at
\EbWpBest\% to win. \EbBreakLabel: \EbBreaks. He played \EbPlayed{} (blue), \EbPlayedSound{} at
\EbWpPlayed\%.}
\label{fig:early}
\end{subfigure}
\caption{Two of his positions with a break on. A break is green where it is sound and red where
it is not; win\% is his after each move. Stockfish at \GradeNodes{} nodes, rerun for this
figure.}
\label{fig:examples}
\end{figure}

\section{What the app can take}

\begin{itemize}[nosep]
\item \textbf{Not the break as the plan whenever it is on.} In his positions a break that can be
  played is sound about one time in five, no more often than a random move, while his own move is
  sound more than half the time. When the app names a structure's break in a position, in the
  review, a drill or the coach's reading of the position, it says whether Stockfish finds the
  break sound there, and does not offer it as the plan otherwise.
\item \textbf{The count of breaks played is no target.} The structure index counts the games
  where he played the break. A break played merely because it can be is mostly unsound, so a
  higher count is not progress by itself; what the index can count instead is the break played
  where it was sound, and the break missed where it was the best.
\item \textbf{Teach the moment.} Where a break is the best move, about one position in ten here,
  he played one about one time in four. Those positions, and the ones where a break was on and
  wrong, are the material for a drill on when to break, each checked by Stockfish and each with
  the structure's trap, which the catalogue already writes. A next round can test whether the
  conditions the app's coaching method asks before a break (pieces pointed at the lines it opens,
  rooks ready for the file, enough force on that side) tell when a break is sound, with its rule
  fixed first.
\end{itemize}

\section{Caveats}

\begin{itemize}[nosep]
\item \textbf{A break stands for a plan.} The break is tested as the next move. The preparation a
  plan asks for first, rooks behind the advance or the trap each structure names, is not tested,
  except as the first chance in a game against the later ones.
\item \textbf{Breaks by square.} The app counts any pawn move to a break's last square as the
  break, a capture landing there included; this study keeps the app's reading.
\item \textbf{The structure of the position.} The structure index names the structure a game
  settles into and holds; this study names it on the board of each position, as the app's
  reading of a position does.
\item \textbf{Sound is a window.} In a quiet position several moves are often within \Window{}
  win\% of the best, so B and A can both be sound; they are paired by position for that reason.
\item \textbf{Many positions a game for him.} Every one of his positions that qualifies enters,
  so a game where a break stays on for many moves counts many times; intervals resample whole
  games for that reason. Strong players give one position a game.
\item \textbf{Mostly the London, and mostly blitz.} About half of his positions are the London
  pyramid; the by-structure table shows the answer is not only the London's. Most strong games
  are blitz; his are rapid.
\item \textbf{Added after the run.} The descriptions of the moment, his own choice, each break
  apart and the strong players' structures were chosen once his halves had been scored. They
  describe; a rule built on them needs positions this study did not see.
\end{itemize}

\end{document}
